\chapter{术语表}\label{app:glossary}

本附录汇总全书使用的关键术语，按拼音字母排序。
释义为检索用的压缩版，不替代正文论证；
“详见”列给出术语首次获得完整定义或分析的章节。
凡属本书自建的教学构件，沿用正文的 \conf{M} 级标注。
英文对照均随词条给出，仅为主要文献用词，不作翻译仲裁。

\begin{longtable}{@{}p{3.0cm}p{9.2cm}p{1.8cm}@{}}
  \toprule
  术语 & 释义 & 详见 \\ \midrule
  \endhead
  \multicolumn{3}{@{}l}{\textbf{A}} \\ \midrule
  AIGC（\en{AI-Generated Content}）&
    人工智能生成内容。第八章把 2026 年奶蛙视频二创的爆发
    归因于以 Seedance 为代表的视频生成模型把生产门槛降到接近零，
    并据此提出“成本不对称”命题。 &
    第八章 \\ \midrule
  \multicolumn{3}{@{}l}{\textbf{B}} \\ \midrule
  半结构化访谈 &
    介于固定问卷与自由谈话之间的访谈形式：按提纲推进、
    不限于既定顺序与措辞。第十一章交付的六题提纲即属此类。 &
    第十一章 \\ \midrule
  编码者一致性 &
    多名编码者对同一材料独立编码时结果的一致程度，
    以 Krippendorff $\alpha$ 或 Cohen's $\kappa$ 报告；
    $\alpha<0.67$ 的特征项应作废重定义，不得带病入表。 &
    第十一章 \\ \midrule
  变异引擎 &
    本书对 AIGC 工具在生产链中所承担环节的称呼：
    以近零成本、高频率产生形象变体，供给第六章所述的
    “内容软”维度。 &
    第八章 \\ \midrule
  变异权 &
    社群对形象进行再创作与变异的实际自由度。
    与法律上的改编权不同：它存在于意义层面，
    不因个体账号被封禁而消失。 &
    第八章 \\ \midrule
  不对称指标 &
    判别设计：功能机制应\keyw{不对称地}改变某指标，
    而消遣 / 算法等竞争解释对其\keyw{对称地}改变；
    两侧表现之差即机制的证据空间。 &
    第七章 \\ \midrule
  不可复核性 &
    关键证据载体已消失或关闭、任何第三方无法再检验的状态。
    奶蛙起源指向的智能体即典型；本书要求在结论处显式登记之。 &
    第四章 \\ \midrule
  \multicolumn{3}{@{}l}{\textbf{C}} \\ \midrule
  采用者类别 &
    Rogers 创新扩散理论中创新者、早期采用者、
    早期 / 后期大多数等类别；第十章用于五校案例的
    形态归位（\conf{M} 级结构假设）。 &
    第十章 \\ \midrule
  参与式文化 &
    受众积极参与生产、分享与改写内容的文化形态；
    本书以之自我定位：教材既是研究也是参与。 &
    第一章 \\ \midrule
  场景语法 &
    衍生语的使用受场景约束的规则系统（“无了”配吃瘪、
    “糖”配反讽）；配错场景即被识别为外行。 &
    第六章 \\ \midrule
  成本不对称 &
    第八章命题：AIGC 使改\keyw{内容}成本近零、
    改\keyw{形式}（音轨、结尾套路）成本仍高，
    由此产生内容多样、形式固化的分化。 &
    第八章 \\ \midrule
  抽象文化 &
    本书工作定义（\conf{M} 级）：以意义悬置、空耳、
    语体戏仿为可观察标记的中文网络风格类型；
    严肃语体承载无厘头内容是其中心装置。 &
    第二章 \\ \midrule
  抽样框 &
    可从中抽取样本的名单或程序。奶蛙受体研究的总体
    不可枚举，故以“配额 + 多入口”替代完整抽样框。 &
    第十一章 \\ \midrule
  错失恐惧（FoMO） &
    对他人正在经历自己缺席的有益经验的普遍性忧虑；
    第十章用于解释新生入学期的同好会纳新增长。 &
    第十章 \\ \midrule
  \multicolumn{3}{@{}l}{\textbf{D}} \\ \midrule
  道德恐慌（\en{moral panic}） &
    Cohen 术语：媒体与权威把低危害的风格现象
    建构为对秩序的威胁的过程；第十章用于警示
    对校园玩梗的运动式清理。 &
    第十章 \\ \midrule
  倒放大法 &
    社群加工手法：把正典音轨倒放，使“我是奶龙，我才是奶龙”
    变为空耳文本；属“捧腹大笑”声音符号的两条加工线之一。 &
    第六章 \\ \midrule
  低强度协商 &
    对立双方以可撤回、不对抗的动作（移除对纪念）
    表达立场并重申边界的互动形态；
    “低强度协商不是冲突的失败，而是冲突的替代”。 &
    第十章 \\ \midrule
  \multicolumn{3}{@{}l}{\textbf{E}} \\ \midrule
  E1/E2/E3 &
    考据伦理三规则：E1 对象限定（只叙述账号、帖子等
    公开对象）；E2 推测隔离（个人境况只说明存在、
    不转述内容）；E3 权利保留（二创形象叙述附署名权归属）。 &
    第四章 \\ \midrule
  二次爆发 &
    SIR 教学模型中新变体输入项触发的第二条感染峰；
    第五章以双峰曲线示意 2024 与 2026 两波热度（\conf{M} 级）。 &
    第五章 \\ \midrule
  \multicolumn{3}{@{}l}{\textbf{F}} \\ \midrule
  反精致 &
    第七章 M3：以故意做差、画质劣化对抗精致内容的姿态；
    其判别问题是“反主流的规范”还是“另一种人设消费”。 &
    第七章 \\ \midrule
  分层装配模型 &
    本书教学综合（\conf{M} 级）：压力层（M1、M2）、
    风格层（M3）与场景层（M4、M5）叠加装配出单一笑点；
    层间可独立证伪。 &
    第七章 \\ \midrule
  复制子（\en{replicator}） &
    Dawkins 术语：能复制自身的单位。第五章把“梗”处理为
    文化复制子，其存续取决于保真度、多产性与存续性。 &
    第五章 \\ \midrule
  \multicolumn{3}{@{}l}{\textbf{G}} \\ \midrule
  关评拉黑 &
    官方奶龙 2026-03 关闭视频评论区并批量拉黑奶蛙账号的
    治理动作（\conf{C} 级）；第九章解读为
    “只能作用于法律层面可执行对象”的不彻底抑制。 &
    第三、九章 \\ \midrule
  \multicolumn{3}{@{}l}{\textbf{H}} \\ \midrule
  换皮奶蛙 / 异形奶蛙 &
    形象谱系两大类：保留奶蛙形体轮廓、仅换配饰与场景者
    为换皮；形体本身偏离（简化、娘化、物化、跨物种）者为异形。 &
    第六章 \\ \midrule
  \multicolumn{3}{@{}l}{\textbf{J}} \\ \midrule
  机械性（\en{mechanical encrustment}） &
    Bergson 笑论：可笑源于“生命的机械化”——
    活人做出僵硬机器般的程式动作；奶蛙的重复仪式形态
    为其当代样本。 &
    第七章 \\ \midrule
  机制矩阵 &
    第七章交付的四列模板（压力源—补偿行为—
    可观察指标—竞争性解释），供书写任何机制假设使用
    （\conf{M} 级）。 &
    第七章 \\ \midrule
  降格（\en{tendentious wit}） &
    Freud 诙谐理论：笑话经由压缩与凝缩绕过心理抑制，
    释放被压抑的攻击或性冲动能量。 &
    第七章 \\ \midrule
  竞争性解释 &
    与目标机制竞争的替代解释（纯消遣、算法推荐、单纯从众）。
    本书纪律：每个机制假设必须并列写出竞争解释与判别设计。 &
    第七章 \\ \midrule
  \multicolumn{3}{@{}l}{\textbf{K}} \\ \midrule
  可传播性（\en{spreadability}） &
    Jenkins 等人术语：内容被受众主动携带、复制、改写的潜力；
    内容的价值边界由受众的可传播性行为划定。 &
    第九章 \\ \midrule
  可传播性单位 &
    本书教学命题：搬运中不依赖原始语境（内容可换、
    署名可弃、载体可换）的文化构件；
    “结尾集体大笑”为奶蛙语料之最（\conf{M} 级）。 &
    第九章 \\ \midrule
  可复现性清单 &
    第十二章交付的十项自查表（引用可追溯、口径并列、
    C 级措辞、E1--E3、M 级自首等）；
    修订者交付新版本前的最低标准（\conf{M} 级）。 &
    第十二章 \\ \midrule
  空耳（\en{soramimi}） &
    把一种语言的语音听成另一种词句的再诠释实践。
    第六章作为社群知识分析，第八章以“声韵母编辑距离”
    为其操作化（\conf{M} 级）。 &
    第六、八章 \\ \midrule
  口径并列 &
    本书证据纪律：同一指标的不同来源数值并列呈现、
    注明来源与日期，不合并、不取均值
    （如 3800 万与 3900 万并存）。 &
    第三章 \\ \midrule
  \multicolumn{3}{@{}l}{\textbf{L}} \\ \midrule
  立牌 &
    奶蛙符号在物理空间的锚点物：把线上迷因固定为
    可合影、可搬运、可巡游的实物，使符号实践进入
    物理空间并接受校方管理。 &
    第十章 \\ \midrule
  流动奶蛙 &
    上海交通大学发起的空间实践：把立牌带到校园各处
    制造偶遇事件。其合影率高于固定装置为可检验推论。 &
    第十章 \\ \midrule
  \multicolumn{3}{@{}l}{\textbf{M}} \\ \midrule
  meme（迷因） &
    Dawkins：通过模仿传播的文化复制子；
    Shifman 将其操作化为内容 / 形式 / 立场三维。 &
    第二、五章 \\ \midrule
  免疫衰减 &
    SIR 教学模型中恢复者随时间回到易感者状态的通道，
    解释同一梗的复发；与二次爆发配套（\conf{M} 级）。 &
    第五章 \\ \midrule
  迷因簇 &
    第五章定义：共享一组可辨识特征的一群变体；
    归并标准 (C1) 特征集、(C2) 谱系、(C3) 识别。 &
    第五章 \\ \midrule
  迷因生态缝隙 &
    控制工具（作用于法律层面）与控制对象
    （部分转移到意义层面）的错位空间；
    “失败的控制不是控制不努力，而是层面错配”。 &
    第九章 \\ \midrule
  模板固化 &
    第八章观察性命题（\conf{M} 级）：AIGC 二创
    在题材上丰富化的同时，形式收敛于少数模板
    （音轨、结尾套路、时长结构）。 &
    第八章 \\ \midrule
  \multicolumn{3}{@{}l}{\textbf{N}} \\ \midrule
  奶蛋 &
    异形奶蛙中简化到极致的产物：形体特征只剩
    “眼睛 + 大笑”，标定谱系的简化极限；
    因“奶蛋轮椅”叙事被赋予强运人设。 &
    第六章 \\ \midrule
  奶龙 &
    第七印象文化传媒 2019 年起开发的动画 IP（正典），
    2020-01-22 于短视频平台上线；
    产业口径见大事年表。 &
    第三章 \\ \midrule
  奶蛙 &
    奶龙二创中诞生的黄色大笑蛙形迷因形象；
    起源指向星野平台某智能体（\conf{C} 级，不可复核）。 &
    第四章 \\ \midrule
  能量节省 &
    Spencer 以来功能主义笑论：笑是心理能量
    经意外通路释放的“节省装置”；第七章古典理论之一。 &
    第七章 \\ \midrule
  逆反认证 &
    第七章 M5：权利人打压反而为迷因提供“官方认证”；
    其判别设计与治理两难分析见第七章。 &
    第七章 \\ \midrule
  \multicolumn{3}{@{}l}{\textbf{P}} \\ \midrule
  判定标准 &
    第八章纪律：任何“是 / 否奶蛙内容”的编码必须有
    事先写定的可执行规则（样本框、检索式、计数规则），
    使他人可重复。 &
    第八章 \\ \midrule
  配额抽样 &
    按关键亚组（年级、院校层次）设下限的非概率抽样；
    无完整抽样框时的替代方案，毕业年级为关键对比组。 &
    第十一章 \\ \midrule
  \multicolumn{3}{@{}l}{\textbf{Q}} \\ \midrule
  情绪代偿 &
    第七章 M1：以“失败—大笑”结构为压力情绪提供
    即时替代性出口；奶蛙语料的中心机制假设（H7.1）。 &
    第七章 \\ \midrule
  群体稀释 &
    第七章 M2：转发把个体压力“稀释”进群体，
    换取“不止我一个人这样”的归属信号；
    与社交笑研究同源。 &
    第七章 \\ \midrule
  \multicolumn{3}{@{}l}{\textbf{R}} \\ \midrule
  $R_0$（基本再生数） &
    SIR 模型中单个感染者在易感者群体中造成的平均新感染数；
    $R_0>1$ 则梗扩散。本书取值为教学示意，非经验估计。 &
    第五章 \\ \midrule
  \multicolumn{3}{@{}l}{\textbf{S}} \\ \midrule
  SIR 型教学模型 &
    第五章把易感—感染—移出流行病模型改造为
    迷因传播的教学模型（\conf{M} 级，非经验模型）。 &
    第五章 \\ \midrule
  三维框架 &
    Shifman 对迷因的操作化：内容（说什么）、
    形式（怎么说）、立场（对谁说、以何姿态说）；
    本书符号分析的工作框架。 &
    第五、六章 \\ \midrule
  社交笑（\en{social laugh}） &
    Provine 笑研究：笑首先是社会性信号
    （群居情境频率约为独处 30 倍），而非仅对幽默的反应。 &
    第七章 \\ \midrule
  社群互惠 &
    第十二章伦理第三线：研究依赖社群材料即负债，
    应以忠实署名、结果回馈、不利用信息差偿还。 &
    第十二章 \\ \midrule
  生产链 &
    第八章整理的奶蛙内容生产环节序列
    （梗图模板 $\to$ 视频剪辑 $\to$ AIGC 生成）；
    各环节成本结构不同，解释变异供给的来源。 &
    第八章 \\ \midrule
  声韵母编辑距离 &
    第八章定义（\conf{M} 级）：按声母 / 韵母 / 声调
    计算两个短语的编辑代价并归一化得 $d$，
    作为空耳成立度的可计算代理指标。 &
    第八章 \\ \midrule
  示例数据 &
    第十一章配套的虚构 40 例 CSV 文件（\conf{M} 级），
    供练习统计流程；任何场合不得引用为代表真实分布。 &
    第十一章 \\ \midrule
  收编（\en{incorporation}） &
    Hebdige：主流与商业把亚文化风格转化为商品、
    消解其抵抗性的过程；第九章指出奶蛙案例中
    收编是双向运动，并提醒先确认抵抗性证据。 &
    第九章 \\ \midrule
  受众外溢 &
    第三章：奶龙受众从亲子向学生群体扩展的三层解释
    （用户年龄结构、二创作者构成、观看场景迁移）。 &
    第三章 \\ \midrule
  双轨分发 &
    正典内容由制作方控制渠道分发，变种内容由
    二创者网络搬运流通；两条轨道的媒介与经济结构均不同。 &
    第二、三章 \\ \midrule
  \multicolumn{3}{@{}l}{\textbf{T}} \\ \midrule
  糖 &
    社群评级词（反讽用法）：指正典内容中幼稚、
    尬感强的特质；社群以讽刺性彩虹屁完成反向评级。 &
    第六章 \\ \midrule
  同好会 &
    校园迷因社群的半正式组织（如北航奶龙同好会）：
    名义功能是水群，实际功能是组织纳新、贺图、
    表情包生产与立牌活动。 &
    第十章 \\ \midrule
  \multicolumn{3}{@{}l}{\textbf{W}} \\ \midrule
  外部性 &
    第八章：二创的生产成本由个体承担、
    拥挤成本（低质内容刷屏）由公共信息流承担；
    平台侧治理机制设计的分析基础。 &
    第八章 \\ \midrule
  未结案证据链 &
    本书对奶蛙起源的定性：指向存在、时间线连贯，
    但关键载体不可复核；只能作教学示范，不作已证事实。 &
    第四章 \\ \midrule
  文本盗猎（\en{textual poaching}） &
    Jenkins（承 de Certeau）：粉丝对文本的
    选择性挪用与再创作；第二章用于定位二创的意义。 &
    第二章 \\ \midrule
  文化资本 &
    Bourdieu 术语：圈内知识（场景语法、谱系识别）
    的持有量构成圈内 / 圈外的边界，是准入测试的给分依据。 &
    第二、五章 \\ \midrule
  误指 &
    考据阶段术语：把错误来源广泛指认为起源
    （如“奶蛙来自某博主”）；误指是考据的起点而非障碍。 &
    第四章 \\ \midrule
  \multicolumn{3}{@{}l}{\textbf{X}} \\ \midrule
  现存样本偏差 &
    以可检索内容为对象的研究只能看到幸存样本
    （未删除、未注销、被记录者）；结论措辞应带
    “在现存材料中”的限定。 &
    第十、十二章 \\ \midrule
  校本化 &
    全国性梗被转写为本校身份标记（“数分奶蛙”
    “航概奶蛙”）；持有并使用即“本校圈内人”的默示声明。 &
    第十章 \\ \midrule
  协商强度谱 &
    第十章教学模型（\conf{M} 级）：按校方反应 /
    学生反应 / 持久物三维给立牌事件分级（0/1/2）；
    北航 9 小时事件合计 2，属低—中强度区间。 &
    第十章 \\ \midrule
  学期节律对齐 &
    第十一章设计支柱：把测量窗口对齐期中与期末考试周，
    使压力机制的效应在数据中可见。 &
    第十一章 \\ \midrule
  行为数据 &
    自愿提供的脱敏聊天记录片段等实际行为痕迹；
    对机制假设的检验优先于自陈问卷。 &
    第十一章 \\ \midrule
  \multicolumn{3}{@{}l}{\textbf{Y}} \\ \midrule
  衍生语 &
    社群从语料中沉淀的固定用语（“无了”“安迪”等），
    受场景语法约束，兼作圈内识别与身份标记。 &
    第六章 \\ \midrule
  引文腐烂（\en{link rot}） &
    网络引文的系统性失效（链接失效、内容被编辑、
    账号注销）；本书三项应对：访问日期、
    现存样本立场、不可复核性登记。 &
    第十二章 \\ \midrule
  一软两硬 &
    第六章对奶蛙符号系统的结构性判断（\conf{M} 级）：
    内容高替换（软）、形式高稳定（硬）、
    立场高约定化（硬）——用最硬的形式承载最软的内容。 &
    第六章 \\ \midrule
  阈限（\en{liminality}） &
    Turner：脱离日常结构、处于角色之间的过渡状态；
    “流动奶蛙合影”被读作微阈限事件（\conf{M} 级诠释）。 &
    第十章 \\ \midrule
  语域错配 &
    本书的命名与体例策略：严肃语域承载无厘头对象；
    其严肃性不能替代证据的严肃性，最终答案见第十二章。 &
    第一、二章 \\ \midrule
  \multicolumn{3}{@{}l}{\textbf{Z}} \\ \midrule
  窄通道规律 &
    第六章命题：变异集中于少数“通道”（道具、配饰、题材），
    形体轮廓等硬核极少被动摇；第八章以成本不对称
    解释其生产端成因。 &
    第六、八章 \\ \midrule
  正典（\en{canon}） &
    权利主体认定的官方设定与内容（正典奶龙）；
    与变种的区分在法律层面与意义层面上并不重合。 &
    第二、三章 \\ \midrule
  证据分级（A/B/C/M） &
    本书制度：A 一手权威、B 机构公开、C 社群考据
    （不作论证前提）、M 本书自建教学构件（须显式声明）；
    区分维度是来源的产生方式，不是内容的真伪。 &
    第一、十二章 \\ \midrule
  证伪条件 &
    第七章纪律：每个机制假设写出“什么观察会使我放弃它”；
    H7.5（逆反认证）的证伪条件设计为范例。 &
    第七章 \\ \midrule
  知情同意 &
    受体研究伦理：施测 / 访谈前告知目的、用途、
    退出权并取得同意；未成年人不在总体内。 &
    第十一章 \\ \midrule
  自陈动机 &
    问卷测得的受访者自述原因；受社会赞许性影响，
    只作弱证据，强检验需行为数据。 &
    第十一章 \\ \midrule
  自我检验 &
    第十二章：本书以自身标准检验自身并列出三处病灶，
    不掩盖、不辩解；方法书最不能做的是
    暗示自己没有病灶。 &
    第十二章 \\ \midrule
  最小化收集 &
    伦理原则：不收集姓名、学号等可识别信息；
    年级用区间，院校由受访者自选披露层级。 &
    第十一章 \\ \bottomrule
\end{longtable}
